\section{Threat model and security property}
\label{sec:threat}

The claims of \cref{sec:resist,sec:prevent} are conditional, and this section
makes the conditions explicit: what is trusted, what the attacker controls, the
property that is meant to hold, and the deployment premises it requires. We state
the property as an invariant rather than rely on the informal ``cannot break
out,'' because the closure theorem the argument inherits~\cite{matchertext} and
the property a security claim actually needs are not the same statement.

\paragraph{Protected asset.}
The host program's parse structure. For a host program with a hole $P[\cdot]$
into which an untrusted value $v$ is placed, the asset is the invariant that the
\emph{structure} of the host's parse of $P[v]$ is independent of the content of
$v$: the value occupies a single data position and contributes no host syntax.

\paragraph{Trusted computing base.}
The guarantee rests on (i) a \emph{matchertext-aware host} that delimits the hole
with a matcher pair and ends the embedded value by matcher balance rather than by
scanning for its native closer; (ii) the \textsc{Verify} recognizer that decides
membership in the matchertext language $L$; and (iii) the total encoder
\textsc{ToMatchertext} and its inverse, used on free-form fields.
\textsc{Verify} is small and host-independent and is the only component the
argument adds to the host, but the boundary-locating step of (i) is the
operationally load-bearing part: we assume it correct here, and note that its
mechanization remains open in the underlying formalism~\cite{matchertext}.

\paragraph{Attacker.}
The attacker supplies $v$ and may choose \emph{arbitrary} bytes, character
encodings, Unicode, length, and matcher-nesting depth, and may submit strings
that are \emph{not} valid matchertext.
Following Kerckhoffs, the attacker knows the discipline, which delimiter kind
guards the hole, and where the hole boundaries lie, and crafts $v$ \emph{against}
the matcher-delimited slot---not the legacy quote- or tag-delimited sink that the
recorded corpus payloads (\cref{sec:results}) target.
The attacker does not control the host source, the configuration $(\Sigma,\Pi)$,
or the trusted components above.

\paragraph{Security property (inertness).}
Let $\textsc{Verify}(v)\Leftrightarrow v\in L$.
For every host program $P[\cdot]$ and every $v$ with $\textsc{Verify}(v)$, the
host parse of $P[v]$ places the content of $v$ at a single terminal whose bytes
do not contribute to the parse structure of $P$; equivalently, the structural
skeleton of the parse of $P[v]$ is independent of $v$.
This is strictly stronger than the closure property $s_o\|v\|s_c\in L$ that
\cref{sec:resist} inherits from the always-embeddability
theorem~\cite{matchertext}: closure states that the \emph{delimited string
remains valid matchertext}, whereas inertness states that the \emph{host parser
treats the delimited content as data}.
Closure is a lemma toward inertness, not inertness itself; bridging the two needs
an explicit host-parser and interpreter model, which we leave to the
implementation work of \cref{sec:concl} and treat as an assumption here.

\paragraph{Deployment premises.}
Inertness holds under four assumptions, each necessary:
\begin{itemize}
\item[\textbf{A1}] \emph{Aware host.} The sink is a matchertext-aware host that
	ends the value by matcher balance rather than at its native closer.
\item[\textbf{A2}] \emph{Delimited context.} Every untrusted value is routed
	into a matcher-delimited hole; a value concatenated into host syntax
	outside a hole is unprotected, exactly as with any point-of-use defense.
\item[\textbf{A3}] \emph{Canonicalize, then verify.} \textsc{Verify} runs on the
	\emph{fully decoded and normalized} value, at the last transform boundary
	before it enters the hole, and the host reads that same representation to
	matcher balance, performing no further decoding, normalization, or
	reparsing of the hole contents.
	Writing $\mathrm{canon}(v)$ for the composition of every decode and
	normalization the pipeline applies, the obligation is
	$\textsc{Verify}(\mathrm{canon}(v))$ with the host parsing
	$\mathrm{canon}(v)$; \textsc{Verify} then need track only the six literal
	matcher code points.
	A check on any earlier, still-encoded representation is unsound: a value
	valid at that point can acquire an unmatched matcher after a later decode
	(\eg \verb|%5D|$\rightarrow$\verb|]|) and break out.\footnote{%
	Making \textsc{Verify} itself robust to \emph{transport} encodings
	(percent, entity, backslash, \dots) is not viable: those forms compose to
	unbounded depth and vary by pipeline, reintroducing the parser-shadowing
	burden matchertext avoids (\cref{sec:bg:sanitizer}).
	Unicode normalization is the one bounded case: the NFKC/NFKD preimage of the
	six matchers---fullwidth, small-form, and presentation variants such as
	\texttt{U+FF3B}$\rightarrow$\texttt{[}---is a finite, UCD-derived set of a
	few dozen code points that \textsc{Verify} may additionally fold or reject
	as a belt-and-suspenders guard against normalization applied after the
	check.}
\item[\textbf{A4}] \emph{Not a semantic sink.} The host does not execute the
	contained value by design (server-side templates, expression languages,
	code evaluation, \verb|javascript:| URLs); delimiting confines structure,
	not meaning.
\end{itemize}

\paragraph{Scope.}
In scope is \emph{structural} injection: an attacker-controlled value that alters
the host parse by escaping its region.
Out of scope, and excluded from the containment claims throughout, are:
execution and semantic sinks (A4); contexts delimited by non-matchers with no
matcher hosting (shell separators, CR/LF header breaks, XML/XXE angle brackets);
\emph{semantic} abuse in which a well-formed value is merely unauthorized, an
access-control question; algorithmic denial of service (regular-expression
backtracking, entity expansion), which is not a breakout; and second-order
injection in which stored input is reinterpreted later in a different context,
unless it is re-checked under A3 at that later sink.
\cref{sec:prevent} measures the in-scope fraction and reports the out-of-scope
residue explicitly.